% GENERATED FILE. Edit canonical lessons, then run npm run generate:books.
% canonical-source-hash: fnv1a64:6898d0d4b35c3e8c
% canonical-lessons: LA-C441-portus, LA-C441-navale, LA-C441-aestuarium, LA-C441-pharus, LA-C441-promunturium

\chapter{Harbour, Dockyard, Tidal inlet, Lighthouse, Headland}
\label{ch:la-harbour-dockyard-tidal-inlet-lighthouse-headland}
\hlchaptermodality{441}

\begin{chapteropening}
\textbf{By the end of this chapter:} \emph{I can describe a harbour and its coast: the harbour, the dockyard, a tidal inlet, the lighthouse and a headland.}

Complete the last lesson of chapter 441: I can describe a harbour and its coast: the harbour, the dockyard, a tidal inlet, the lighthouse and a headland.
\end{chapteropening}

\section[portus, portūs]{portus, portūs — a harbour, a port}

\label{lesson:LA-C441-portus}



\begin{quote}
Before the new one: say the Latin for to take across, then the Latin for to swim to safety.
\end{quote}

\subsection*{You'll want to know: portus, portūs}
\textbf{portus, portūs} — \textquotedblleft{}a harbour, a port\textquotedblright{}.

\begin{quote}
Tandem nāvis in portum intrat, et nautae gaudent.
\end{quote}

\begin{quote}
\emph{At last the ship comes into the harbour, and the sailors are glad.}
\end{quote}

Fourth declension, masculine: \emph{portus, portūs m.} Figurative too: a safe place after trouble. English \emph{port} comes from it.

\begin{cousinweb}
\emph{portus} is inherited from Proto-Italic \emph{portus}, from PIE \emph{pértus} 'crossing, passage' (EtymDB), the same root as \emph{porta} 'gate'. English \emph{port} comes from it.
\end{cousinweb}

\begin{grammarlens}[title={What you've built}]
One more word for telling, explaining and describing at B1.
\end{grammarlens}

\subsection*{Your turn}
\begin{itemize}
\raggedright
  \item \emph{Say it:} \emph{portus, portūs}
  \item \emph{Say it:} \emph{portus, portūs}, once more
  \item \emph{Say it:} a sentence of your own with \emph{portus, portūs}
  \item \emph{Recall:} say the Latin for to take across, then the Latin for to swim to safety, then say \emph{portus, portūs} again
\end{itemize}

\begin{tcolorbox}[breakable,colback=teal!4,colframe=teal!35!black,title={Before you move on}]
What does portus, portūs mean? (\textquotedblleft{}A harbour, a port\textquotedblright{}.) Say it once more.
\end{tcolorbox}
\section[nāvāle, nāvālis]{nāvāle, nāvālis — a dockyard, a place where ships are kept}

\label{lesson:LA-C441-navale}



\begin{quote}
Before the new one: say the Latin for to swim to safety, then the Latin for a harbour.
\end{quote}

\subsection*{You'll want to know: nāvāle, nāvālis}
\textbf{nāvāle, nāvālis} — \textquotedblleft{}a dockyard, a place where ships are kept\textquotedblright{}.

\begin{quote}
In nāvālī fabrī nāvēs novās aedificant.
\end{quote}

\begin{quote}
\emph{In the dockyard the workmen build new ships.}
\end{quote}

Third declension, neuter: \emph{nāvāle, nāvālis n.}, often plural \emph{nāvālia}. It is the neuter of the adjective \emph{nāvālis} 'of ships'; English \emph{naval} comes from that.

\begin{cousinweb}
\emph{nāvāle} is the neuter of the adjective \emph{nāvālis} 'of ships', built on \emph{nāvis} 'ship' (Lewis \& Short). English \emph{naval} comes from the adjective.
\end{cousinweb}

\begin{grammarlens}[title={What you've built}]
One more word for telling, explaining and describing at B1.
\end{grammarlens}

\subsection*{Your turn}
\begin{itemize}
\raggedright
  \item \emph{Say it:} \emph{nāvāle, nāvālis}
  \item \emph{Say it:} \emph{nāvāle, nāvālis}, once more
  \item \emph{Say it:} a sentence of your own with \emph{nāvāle, nāvālis}
  \item \emph{Recall:} say the Latin for to swim to safety, then the Latin for a harbour, then say \emph{nāvāle, nāvālis} again
\end{itemize}

\begin{tcolorbox}[breakable,colback=teal!4,colframe=teal!35!black,title={Before you move on}]
What does nāvāle, nāvālis mean? (\textquotedblleft{}A dockyard, a place where ships are kept\textquotedblright{}.) Say it once more.
\end{tcolorbox}
\section[aestuārium, aestuāriī]{aestuārium, aestuāriī — a tidal inlet, a tidal marsh; an estuary}

\label{lesson:LA-C441-aestuarium}



\begin{quote}
Before the new one: say the Latin for a harbour, then the Latin for a dockyard.
\end{quote}

\subsection*{You'll want to know: aestuārium, aestuāriī}
\textbf{aestuārium, aestuāriī} — \textquotedblleft{}a tidal inlet, a tidal marsh; an estuary\textquotedblright{}.

\begin{quote}
Ubi mare recēdit, in aestuāriō lutum est.
\end{quote}

\begin{quote}
\emph{When the sea goes back, there is mud in the tidal inlet.}
\end{quote}

Second declension, neuter: \emph{aestuārium, aestuāriī n.} Land the tide covers and uncovers. English \emph{estuary} comes from it.

\begin{cousinweb}
\emph{aestuārium} is built on \emph{aestus} 'heat, surge, tide' with the suffix \emph{-ārium} (EtymDB; Lewis \& Short). English \emph{estuary} comes from it.
\end{cousinweb}

\begin{grammarlens}[title={What you've built}]
One more word for telling, explaining and describing at B1.
\end{grammarlens}

\subsection*{Your turn}
\begin{itemize}
\raggedright
  \item \emph{Say it:} \emph{aestuārium, aestuāriī}
  \item \emph{Say it:} \emph{aestuārium, aestuāriī}, once more
  \item \emph{Say it:} a sentence of your own with \emph{aestuārium, aestuāriī}
  \item \emph{Recall:} say the Latin for a harbour, then the Latin for a dockyard, then say \emph{aestuārium, aestuāriī} again
\end{itemize}

\begin{tcolorbox}[breakable,colback=teal!4,colframe=teal!35!black,title={Before you move on}]
What does aestuārium, aestuāriī mean? (\textquotedblleft{}A tidal inlet, a tidal marsh; an estuary\textquotedblright{}.) Say it once more.
\end{tcolorbox}
\section[pharus, pharī]{pharus, pharī — a lighthouse}

\label{lesson:LA-C441-pharus}



\begin{quote}
Before the new one: say the Latin for a dockyard, then the Latin for a tidal inlet.
\end{quote}

\subsection*{You'll want to know: pharus, pharī}
\textbf{pharus, pharī} — \textquotedblleft{}a lighthouse\textquotedblright{}.

\begin{quote}
Nocte lūx pharī procul ardet, et nautae portum inveniunt.
\end{quote}

\begin{quote}
\emph{At night the light of the lighthouse burns far off, and the sailors find the harbour.}
\end{quote}

Second declension, feminine: \emph{pharus, pharī f.} (also \emph{pharos}). Named after a famous lighthouse island.

\begin{cousinweb}
\emph{pharus} was borrowed from Greek \emph{pharos}, the name of the island whose lighthouse became famous (EtymDB; Lewis \& Short); the name then became the common word for any lighthouse.
\end{cousinweb}

\begin{grammarlens}[title={What you've built}]
One more word for telling, explaining and describing at B1.
\end{grammarlens}

\subsection*{Your turn}
\begin{itemize}
\raggedright
  \item \emph{Say it:} \emph{pharus, pharī}
  \item \emph{Say it:} \emph{pharus, pharī}, once more
  \item \emph{Say it:} a sentence of your own with \emph{pharus, pharī}
  \item \emph{Recall:} say the Latin for a dockyard, then the Latin for a tidal inlet, then say \emph{pharus, pharī} again
\end{itemize}

\begin{tcolorbox}[breakable,colback=teal!4,colframe=teal!35!black,title={Before you move on}]
What does pharus, pharī mean? (\textquotedblleft{}A lighthouse\textquotedblright{}.) Say it once more.
\end{tcolorbox}
\section[prōmunturium, prōmunturiī]{prōmunturium, prōmunturiī — a headland, a promontory}

\label{lesson:LA-C441-promunturium}



\begin{quote}
Before the new one: say the Latin for a tidal inlet, then the Latin for a lighthouse.
\end{quote}

\subsection*{You'll want to know: prōmunturium, prōmunturiī}
\textbf{prōmunturium, prōmunturiī} — \textquotedblleft{}a headland, a promontory\textquotedblright{}.

\begin{quote}
Nāvis prope prōmunturium altum nāvigat.
\end{quote}

\begin{quote}
\emph{The ship sails near a high headland.}
\end{quote}

Second declension, neuter: \emph{prōmunturium, prōmunturiī n.}, also spelled \emph{prōmontōrium}. English \emph{promontory} comes from it.

\begin{cousinweb}
\emph{prōmunturium} is usually traced to \emph{prō} 'forward' and \emph{mineō} 'to jut out' (Lewis \& Short): land that juts forward into the sea. English \emph{promontory} comes from it.
\end{cousinweb}

\begin{grammarlens}[title={What you've built}]
One more word for telling, explaining and describing at B1.
\end{grammarlens}

\subsection*{Your turn}
\begin{itemize}
\raggedright
  \item \emph{Say it:} \emph{prōmunturium, prōmunturiī}
  \item \emph{Say it:} \emph{prōmunturium, prōmunturiī}, once more
  \item \emph{Say it:} a sentence of your own with \emph{prōmunturium, prōmunturiī}
  \item \emph{Recall:} say the Latin for a tidal inlet, then the Latin for a lighthouse, then say \emph{prōmunturium, prōmunturiī} again
\end{itemize}

\begin{tcolorbox}[breakable,colback=teal!4,colframe=teal!35!black,title={Before you move on}]
What does prōmunturium, prōmunturiī mean? (\textquotedblleft{}A headland, a promontory\textquotedblright{}.) Say it once more.
\end{tcolorbox}
